\section{Effective Actions and their 12d Consistencies}
The type IIB supergravity \eqref{Einstein frame IIB action} is understood as the low energy effective field theory of the type IIB string theory.
Under higher-derivative and stringy corrections, the SL(2, R) symmetry is believed to be broken to the discrete group SL(2, Z) \cite{Schwarz:1995dk}.
If the type IIB theory admits certain 12d interpretation, such interpretation, and its implications on effective actions, must be consistent with the type IIB effective action and SL(2, Z).

In this section, we review the current understandings in SL(2, Z) as arising from higher-derivative corrections, and their 12d interpretations.
We will begin by reviewing how one obtains higher-derivative corrections in supergravity from perturbative string theory, most importantly how the SL(2, Z) invariance arises non-perturbatively from the D(-1) backgrounds.
Then we discuss current understandings on 12d interpretations of SL(2, Z), and identify certain directions in advancing them. 
We conclude by discussing a potential parallel, between type IIA and type IIB, where the D0 and D(-1) backgrounds are effectively accounted for by loops of higher-dimensional KK modes.

\subsection{Recap: Perturbative String Theory and Higher-Derivative Supergravity}
{\color{orange}
The perturbative expansion of string amplitudes is organized by two parameters: $\alpha' = l_s^2$ controls the derivative expansion, and $g_s$ counts the genus of the string worldsheet in string path-integrals.
An $n$-point amplitude is a summation over worldsheet path-integrals of genus $g$
\begin{equation}
    \mathcal A_n(\alpha') = \sum_{g=0}^\infty g_s^{2g-2+n} \mathcal A_n^{(g)}(\alpha').
\end{equation}
}


The tree-level 4-point function in the type IIA and type IIB string theories is the Virasoro-Shapiro amplitude \cite{Green:2012oqa}, for the symmetric traceless massless modes in the NSNS sector of the string, they take the following form \cite{Green:1999pv} up to an irrelevant overall factor
\begin{equation}
\begin{aligned}
\mathcal A_4
&=\frac{\alpha^{\prime 4}}{3\cdot 2^3 g_s^2}
K
\times\frac{64}{\alpha^{\prime 3}stu}
\frac{\Gamma[1-\tfrac{\alpha'}{4}s]\Gamma[1-\tfrac{\alpha'}{4}t]\Gamma[1-\tfrac{\alpha'}{4}u]}{\Gamma[1+\tfrac{\alpha'}{4}s]\Gamma[1+\tfrac{\alpha'}{4}t]\Gamma[1+\tfrac{\alpha'}{4}u]}\\
K &\equiv t_8^{m_1...m_8}t_8^{n_1...n_8}               
\prod_{r=1}^4\zeta^{(r)}_{m_{2r}n_{2r}}k^{(r)}_{m_{2r-1}}k^{(r)}_{n_{2r-1}}
\end{aligned}
\end{equation}
where $s,t,u$ are Mandelstam variables, $\zeta^{(i)}_{mn}$ is the dimensionless polarization of the symmetric traceless modes on a closed string, they may also be interpreted as polarizations of the spacetime graviton $\kappa h^{(i)}_{mn}$ in $G_{mn}=\eta_{mn} + \kappa h_{mn}$. 
The explicit form of $t_8 t_8$ can be found in \cite{Schwarz:1982jn,Policastro:2008hg,Liu:2025uqu}.
Using the relation
\begin{equation}
\ln\Gamma(1-x)-\ln\Gamma(1+x) = 2\gamma x + 2\sum_{m\ge0}\frac{\zeta(2m+3)}{2m+3}x^{2m+3},
\end{equation}
we find the expansion
\begin{equation}
\mathcal{A}_4 = \frac{\alpha^{\prime 4}}{3\cdot 2^3 g_s^2}
K
\left[ \frac{64}{\alpha^{\prime 3}stu} + 2 \zeta(3)\right]
 + O(\alpha^{\prime 5})
\end{equation}
evaluated in the string frame metric with $G_{mn} |_{\infty} = \eta_{mn}$.
From the relation $G_{mn} = e^{\Phi/2}g_{mn}$, we see that the Einstein frame metric is not asymptotically flat: $g_{mn} |_{\infty} = g_s^{1/2}\eta_{mn}$.
To translate $\mathcal{A}_4$ into the Einstein frame, we need to translate implicit contractions with the string frame metric with those with the Einstein frame metric, which introduces factors of $g_s$.
The net effect is the following Einstein frame expansion
\begin{equation}
\mathcal{A}_4 \bigg|_{g_{mn}} = 
\frac{\alpha^{\prime 4}}{3\cdot 2^3 }
K
\times
\left[ \frac{64}{\alpha^{\prime 3}stu} + 2 \zeta(3)g_s^{-3/2}\right]
 + O(\alpha^{\prime 5}),
\end{equation}
with $\zeta_{mn}^{(r)}$ inside $K$ now interpreted as the polarization of the Einstein frame graviton and all contractions are implicitly performed with the asymptotic Einstein frame metric $g_s^{1/2}\eta_{mn}$.

The polynomial $K$ is analytic in the momenta and is the linear term in the expansion of $\frac{t_8 t_8 R^4}{16}$.
The nonlocal poles at $s, t, u = 0$ correspond to massless graviton exchange and reproduce the $s,t,u$-channel structure generated by the cubic graviton vertices of the Einstein-Hilbert action; the nonlocal $1/stu$ piece is therefore attributed to the tree-level supergravity dynamics.
Indeed it is accompanied by a single factor of $\alpha'$ which reveals its origin in the two-derivative supergravity.

The next contribution appears at $O(\alpha^{\prime 4}$) which reveals an 8-derivative spacetime origin:
\begin{equation}
\frac{\alpha^{\prime 4}}{3 \cdot 2^3} K 2 \zeta(3) g_s^{-3/2} \cong
\frac{\alpha^{\prime 4}}{3 \cdot 2^7} t_8t_8 R^4 (2 \zeta(3) g_s^{-3/2}) + O(\alpha^{\prime 5}).
\end{equation}
The next order correction comes from genus-1 amplitudes \cite{Green:1999pv}.
It is also common in both type IIA and IIB theories, and comes with the $t_8 t_8 R^4$ factor. 
Combining the genus-0 and genus-1 amplitudes, we have the local effective actions
\begin{equation}
\mathcal{A}_4 = -4\pi^7 (t_8 t_8 R^4)(\alpha')^4
\left[
2 \zeta(3)\tau_2^{3/2}
+\frac{2\pi^2}{3}\tau_2^{-1/2}
\right]
+O(\alpha^{\prime 5}).
\end{equation}
Non-renormalization theorems suggest the perturbative corrections to $R^4$ terminate here at one-loop \cite{Green:1997tv,Kehagias:1997cq}.
Note that $\mathcal{A}_4$ as given above no longer has SL(2, R) symmetry. 
In fact, the SL(2, R) symmetry of the type IIB supergravity is only present at the two-derivative level together with corrections that do not depend on $\tau$.
As soon as higher-derivative terms with non-trivial dependence on $\tau$ enter, SL(2, R) is explicitly broken.

In particular, the type IIB string path integral requires summing over the D(-1) backgrounds \eqref{D-1 background}.
The single and multi-charged D(-1) backgrounds give rise to non-perturbative $R^4$ corrections \cite{Green:1997tv,Kehagias:1997cq} of the form
\begin{equation}
\mathcal A_{D(-1)}\propto (\alpha')^4t_8t_8R^4\sum_{m,n\ge1}\left(\frac{m}{n^3}\right)^{1/2}
(e^{2\pi i mn\tau}+e^{-2\pi i mn\bar{\tau}})
\left(
1+\sum_{k=1}^\infty(4\pi mn\tau_2)^{-k}\frac{\Gamma[k-1/2]}{\Gamma[-k-1/2]k!}
\right).
\end{equation}
When combined with the genus-0 and genus-1 perturbative contributions, these D(-1) terms assemble into the modular-invariant, non-holomorphic Eisenstein series
\begin{equation}
\begin{aligned}
f_0(\tau,\bar\tau)
&=\sum_{(m,n)\neq(0,0)}\frac{\tau_2^{3/2}}{|m+n\tau|^3}\\
&=2\zeta(3)\tau_2^{3/2}+\frac{2\pi^2}{3}\tau_2^{-1/2}\\
&\quad+4\pi^{3/2}\sum_{m,n\ge1}\left(\frac{m}{n^3}\right)^{1/2}
(e^{2\pi i mn\tau}+e^{-2\pi i mn\bar{\tau}})
\left(
1+\sum_{k=1}^\infty(4\pi mn\tau_2)^{-k}\frac{\Gamma[k-1/2]}{\Gamma[-k-1/2]k!}
\right),
\end{aligned}
\label{eq:f_0 Eisenstein}
\end{equation}
thereby restoring the SL(2, Z) duality symmetry.

\subsection{12d Effective Corrections}
\label{sec:12d effective corrections}
From our previous discussion, the type IIB 4-graviton effective action takes on the following schematic form \cite{Kehagias:1997cq,Green:1997tv}
\begin{equation}
L^{(3)} \propto f_0(\tau,\bar\tau)t_8t_8 R^4
\label{10d IIB 4-graviton effective action}
\end{equation}
The rest of the 4-point effective action consists of axio-dilaton and is SL(2, Z) invariant, they can be found in \cite{Policastro:2006vt,Policastro:2008hg}
To produce the 10d effective action \eqref{10d IIB 4-graviton effective action},
a ``12d effective action" had been proposed as \cite{Minasian:2015bxa} 
\begin{equation}
\hat{L}^{(3)} \propto f_0(\tau,\bar\tau)t_8t_8 \hat{R}^4,
\label{eq:liu-12d-ansatze}
\end{equation}
where $\hat{R}$ is the 12d Riemann tensor computed with $\hat{g}_{MN}$, as in \eqref{12d metric ansatz}
It was shown that \eqref{eq:liu-12d-ansatze} reduced on the 12d metric ansatz  produces, at 4-point, the effective action in the axio-dilaton sector \cite{Policastro:2006vt,Policastro:2008hg}.
However, it was later shown that \eqref{eq:liu-12d-ansatze} is inconsistent with 10d type IIB amplitudes at 5-point \cite{Liu:2025uqu}.

We now comment on certain limitations of the proposed 12d effective action \eqref{eq:liu-12d-ansatze} and outline possible directions forward.
First, $\epsilon_{10}\epsilon_{10}R^4$ vanishes at 4-point, so the result of \cite{Minasian:2015bxa} on $\epsilon_{10}\epsilon_{10}R^4$ was that $\varepsilon_{12}\varepsilon_{12}\mathcal R^4$ vanishes at 4-point as well, after reducing $\mathcal G_{MN}$ to $g_{mn},\phi,C$.
The correspondence would be significantly stronger if non-vanishing components of $\epsilon_{10}\epsilon_{10}R^4$ could be verified, e.g. for 5-point amplitudes.
Unfortunately, this does not occur \cite{Liu:2025uqu}.

A second issue concerns the interpretation of the $t_8t_8 R^4$ term.
The $t_8$ tensor is originally defined by traces over the gamma matrices of SO(8) \cite{Schwarz:1982jn}, which is the little group of SO(9, 1).
Accordingly, the kinematic structure encoded by $t_8t_8R^4$ is that of the 8d space transverse to a massless momentum.
In the 4-graviton amplitude, the kinematic structure of $t_8t_8 R^4$ is thus determined solely by the 8 transverse components of the graviton polarizations and momenta, rather than all 10.
For example, one is able to extract $t_8t_8 R^4$ from 9d amplitudes obtained by compactifying M-theory on a torus \cite{Green:1997as}.
Consequently, when examining non-vanishing 4-graviton $t_8t_8R^4$ amplitudes, it is ambiguous whether one is investigating the established relation between 11d and 9d, or between 12d and 10d.
By contrast, $\epsilon_{10}\epsilon_{10}R^4$ is intrinsically 10d.
Thus, to strengthen the proposed relation between 12d and 10d amplitudes, it is worth investigating how one may capture $\epsilon_{10}\epsilon_{10}R^4$ from 12d.

Lastly, from a geometric perspective,in 12d $\tau$ should not explicitly appear. 
The point in repackaging the type IIB theory and its corrections in a 12d-covariant way is to geometrize $\tau$ as part of the metric, so one should be alarmed to find the need to put in SL(2, Z) covariance by hand, e.g. via $f_0(\tau,\bar\tau)$.
Perhaps a more appropriate 12d amplitude would be
\begin{equation}
\hat{L}^{(3)} \propto f(\hat{g}_{MN},\hat{R}_{MNPQ}).
\end{equation}
Then upon reduction on a torus, 2 of the 12 directions are singled out, we are thus able to distinguish $\tau$ from the rest of the metric, and obtain the modular function and $R^4$
\begin{equation}
f(\hat{g}_{MN},\hat{R}_{MNPQ})
= f_0(\tau,\bar\tau)(t_8t_8 + \epsilon_{10}\epsilon_{10})R^4 + ....
\end{equation}
This alternative route may be worth exploring, but at the moment it is not clear how to construct such $f(\hat{g},\hat{R})$, and the ansatze of \cite{Minasian:2015bxa} remains the most natural choice.
One may look into functions that admit expansions over $f_0$, or consider possible 12d interpretations of the IIB 5-point amplitudes.
The type IIB amplitudes, starting at 5-points, famously contain the ``U(1)-violating terms" \cite{Liu:2025uqu}. This is one clear obstruction to extending the 12d ansatze of \eqref{eq:liu-12d-ansatze}.
It would also be illuminating to elucidate how this obstruction shall be interpreted, or worked around, in 12d.
Alternatively, one may still hope that certain 12d tensor contraction structure could capture the kinematics of 10d amplitudes, as constrained by SL(2, Z), and be content with a separate ansatze for the U(1)-violating sectors.





\subsection{KK-modes and D-brane Backgrounds}
\label{sec:KK modes}
The 12d perspective suggests one further parallel between the type IIA and type IIB effective actions, concerning the role of supergravity KK-modes as surrogates for the D0 and D(-1) backgrounds in the string path integral.
In the type IIA theory, the tower of KK-modes on the M-theory circle $S_1$, running in a loop, reproduces the $t_8t_8R^4$ coupling together with its D0-brane (worldline instanton) completion \cite{Green:1997as}.
The analogue on the type IIB side is that the KK-modes of a 12d field on $T_2$ reproduce the genus-1 and D(-1) contributions to the $R^4$ coupling.

Compactifying $(x^\mu,y^1,y^2)$ on a $T_2$ of complex structure $\tau$ and radius $R$, a 12d scalar acquires the KK spectrum $M_{m,n}^2=|m+n\tau|^2/(R^2\tau_2^2)$.
Evaluating the one-loop four-graviton box with this tower running inside --- via the Schwinger representation, performing the Poisson resummation and zeta-function renormalization explicitly (Appendix \ref{sec:KK loop R4}), or by the shortcut of adding a spectator dimension to the known 11d-on-$T_2$ results \cite{Green:1997as,Beccaria:2023hhi} --- one finds
\begin{equation}
\mathcal{A}_4
\propto  f_0(\tau,\bar\tau)(s^2+t^2+u^2)+\ldots
\sim f_0(\tau,\bar\tau)t_8t_8R^4+\ldots.
\end{equation}
The single KK-loop sum thus reproduces the entire modular coupling $f_0$ of \eqref{eq:f_0 Eisenstein}: the tree and genus-1 terms appear as its constant part, and the D(-1) instanton series as its exponentially suppressed modes.
In this sense the KK-tower on $T_2$ acts as a surrogate for the D(-1) background, in parallel with the type IIA KK-tower on $S_1$ and the D0 background \cite{Green:1997as}.

We do not wish to overclaim.
As shown in Appendix \ref{sec:KK loop R4}, the modular weight of $f_0$ is fixed by the proper-time power $\lambda^{-3/2}$, which is the one-loop box measure in \emph{nine} noncompact dimensions; the computation is therefore the established 11d-on-$T_2$ relation \cite{Green:1997as} carried along by a spectator tenth dimension, rather than an intrinsically 12d-to-10d statement.
This is consistent with $t_8t_8R^4$ not being an intrinsically 10d term, as discussed in Section \ref{sec:12d effective corrections}.
A genuine 12d-to-10d diagnostic would require an amplitude that signals all ten momenta, e.g. $\epsilon_{10}\epsilon_{10}R^4$, which first appears at 5-point.
One may attempt the analogous KK-loop evaluation at 5-point using the proper-time formulae of \cite{Minasian:2015bxa,Green:1999by,Liu:2022bfg}.